\cotefile{103}
\scanpage{11}
\deleted{\unsure{catégories|catégorie}}\inserted{\unsure{champs}} de Picard \unsure{épinglés} par $M,N$, $\unsure{\sigma|o}$ étant l'invariant qui \inserted{s'obtient en} associant à toute section \inserted{$L$} d'un champ de Picard la symétrie de $\deleted{\unsure{K|L}}\,L\otimes L$, interprétée comme section de ${}_{2}N$. Je sais prouver (sauf erreur) que tout hom. $M\to{}_{2}N$ provient d'un champ de Picard convenable (épinglé par $M,N$) (\unsure{a|à} priori l'obstruction est dans $\operatorname{Ext}^{3}(X;M,N)$, mais un argument ``universel'' prouve qu'elle est nulle). Cela prouve que l'extension $\alpha$, est bien proche d'être splittée :\deleted{\unsure{u|un}} tte section du troisième faisceau, sur un objet quelconque de $X$, se remonte - en d'autres termes l'extension a une section ``ensembliste''. Bien sûr, il y a mieux en fait : toute section sur un $U\in\operatorname{Ob}X$ ``provient'' d'un élément de ${\unsure{\mathcal H|\mathscr H}}^{2}(U,\unsure{E'|E^{1}}(M,N))$\unsure{,|.} (hypercohomologie $\mathcal H^{2}$).
% archive sheet number: 10
\scanpage{12}

\uline{Exemple} Soit $A$ un anneau sur $X$, soient $M,N$ respectivement les faisceaux $K^{0}$ et $K^{1}$ associés au champ additif des $A$-Modules projectifs de $\unsure{\mathrm{v}|\mathrm{t}}.\mathrm{f}$. \inserted{(p\textsuperscript{r} ex.)} Alors la construction de Mme Sinh nous fournit un champ de Picard épinglé par $M,N$, d'où une section \deleted{\unsure{du|de}} canonique du terme médian $\unsure{F|P}(M,N)$ de $(\alpha)$.

\uline{NB} Tout ce qui précède a les fonctorialités évidentes en $M,N,X$\ldots

\uline{Question} : Le triangle exact $(\unsure{T|J})$ et \unsure{-} la suite exacte \unsure{$\alpha$)|($\alpha$)|$\alpha$} sont ils connus par les compétences (Quillen, Breen, Illusie\ldots) ? \deleted{\unsure{S|X}}Connaissent-ils des variantes ``supérieures'' ? (Un principe ``géométrique'' pour en obtenir pourrait être via des $n$-champs de Picard non néc. stricts\ldots)
% archive sheet number: 11
\scanpage{13}

Je profite de l'occasion pour soulever une question sur la ``cohomologie relative''. Soit $g:X\to Y$ un morphisme de topos. Si $F$ est un faisceau abélien (ou un complexe d'iceux) sur $Y$, peut-on définir \uline{fonctoriellement} en $F$ la cohomologie relative $\mathbb R\Gamma(Y\mathbin{\mathrm{mod}}X,F)$ (de la catégorie dérivée de $\mathit{Ab}(Y)$ dans celle de $\mathit{Ab}$) ? L'interprétation ``géométrique'' en termes d'opérations sur des $n$-champs de Picard ($n$ ``grand'') suggère que ça doit exister. Mais je ne vois de construction \unsure{évidente} ``à la main'' que dans les deux cas extrêmes :

a) $g$ est ``$(-1)$-acyclique'' i.e. pour tt $F$ sur $Y$, $F\unsure{\mathord{-}|\to}g_{*}g^{*}F$ est \unsure{injectif} (\uline{NB} c'est le cas de $Y_{/p}\to Y$ si $p\to e_{Y}$ est un \unsure{épimorphisme} - c'est donc le cas de $B_{\unsure{e|c}}\to B_{\unsure{G|c}}$ plus haut)

On prend
\[
\mathbb R\Gamma_{\deleted{\unsure{Y|X}}}\bigl(\operatorname{Coker}\bigl(F\to g_{*}\unsure{[|(}\underbrace{e(g^{*}(F}_{\text{résol. inj.}})))\bigr)[-1]\bigr)
\]
% archive sheet number: 12
\scanpage{14}

b) $\forall\,F$ injectif sur $\unsure{Y}$, $g^{*}(F)$ est injectif et $F\unsure{\mathord{-}|\to}g_{*}g^{*}F$ est un épimorphisme (ex: $g$ inclusion d'un ouvert $U\hookrightarrow e_{Y}$).

On prend
\[
\mathbb R\Gamma_{Y}\bigl(\operatorname{Ker}\bigl(\underbrace{C(F}_{\text{rés. inj.}})\to g_{*}g^{*}\bigl(C(F)\bigr)\bigr)\bigr)
\]

Dans le cas général, la difficulté provient du fait que le cône d'un morphisme de complexes (tel que
\[
F\to g_{*}\deleted{\unsure{C|(}}\bigl(g^{*}(F)\deleted{\unsure{)|C}}\bigr)\deleted{\unsure{g_{*}|*}}\qquad\text{)}
\]
n'est pas fonctoriel (dans la catégorie dérivée) p. r. à la flèche dont on veut prendre le cône.

Et pourtant, dans le cas particulier actuel il devrait y avoir un choix fonctoriel. \unsure{ESt|Est}-ce \unsure{évident|evident} ?

\uline{Question pour Illusie} : Dans sa théorie des \unsure{déformations|déform.tions} de schémas en \unsure{groupes|groupe} plats, il tombe sur des $H^{3}(B_{G}/X,\quad)$ resp. des $\operatorname{Ext}^{2}(X;,\quad)$. Peut-on courtcircuiter sa théorie via la théorie (supposée écrite) des \unsure{Gr|gr}-champs -- resp. via la théorie des champs de Picard ? J'ai
% archive sheet number: 13
\scanpage{15}

Je te signale que j'ai \unsure{réfléchi} aux Gr-champs sur $X$. Si $G$ est un Groupe sur $X$, $N$ un $G$-Module les Gr-champs sur $X$ ``épinglés par $G,N$'' forment a priori une $2$-\unsure{catégorie} et même une $2$-catégorie de Picard stricte, grâce à l'opération \unsure{évident|évidente} à la Baer. On trouve que le complexe \inserted{(de \unsure{cochaines|cochaînes})} tronqué à l'échelon $2$ qui leur correspond est le tronqué
\[
\tau_{\leq 2}\bigl(\mathbb R\Gamma(B_{G}\mathbin{\mathrm{mod}}X,N)[1]\bigr)
\]
(\uline{NB} la cohomologie de $\mathbb R\Gamma(B_{G}\mathbin{\mathrm{mod}}X,N)$ commence en degré $1$) Plus \unsure{géométriquement}, un Gr-champ sur $X$ \unsure{épinglé} \inserted{par $(G,N)$} est essentiellement ``la même chose'' qu'une $2$-gerbe sur $B_{G}$, liée par $N$, et \deleted{\unsure{pri|épi}}\unsure{munie} d'une \unsure{trivialisation} au dessus de $X\approx B_{\unsure{e|c}}=(B_{G})_{/P}$ (où $P$ est le \deleted{\unsure{s|$\ast$}}objet de $B_{G}$ ``torseur universel sous $G$''). Ces $2$-gerbes forment en fait une $3$-catégorie de Picard a priori, mais il se
% archive sheet number: 14
\scanpage{16}
trouve que dans celle-ci, les $3$-flèches sont triviales (i.e. si source $\unsure{=|\simeq|\ne}$ but, ce sont des identités) - cela un fait qui exprime $H^{0}(B_{G}/X,N)=0$ (i.e. $H^{0}(B_{G},N)\to H^{0}(X,N)$ injectif\ldots). Donc la $3$-catégorie peut être \unsure{regardée|regardé} comme une $2$-catégorie - et ``c'est'' celle des Gr-champs sur $X$ épinglés par $G,N$. Si on veut localiser sur $X$, et décrire le $2$-\uline{champ} \inserted{de Picard} sur $X$ des champs de Picard (sur des objets variables de $X$) épinglés par $G,N$, on trouve qu'il est \unsure{exprime|exprimé} par le complexe
\[
\tau_{\leq 2}\bigl(\mathbb R p_{G*}\operatorname{Coker}\bigl(N\to Rq_{G*}\overbrace{\unsure{C|e}(q_{G}^{*}N)}^{\text{rés. inj.}}\bigr)\bigr)
\]
où $p_{G}:B_{G}\to X$ et $q_{G}:B_{e}\approx X\simeq(B_{G})_{/\unsure{p|P}}\to B_{G}$

Toutes ces descriptions étant compatibles avec des variations de $G,N,X$, cela \unsure{do|de}\deleted{\unsure{nne|une|n}} en principe une description de la $2$-catégorie des Gr-champs, avec $X,G,N$ variables\ldots
% archive sheet number: 15
\scanpage{17}

\noindent\hfill\unsure{Villecun|Villecui} le 23.6.1974\\

Chers Deligne, Giraud, Verdier,\\

Vous savez peut-être qu'une mathématicienne vietnamienne, Hoang Xuan Sinh, \inserted{a} travaillé depuis mon séjour au Vietnam en 1967 à un sujet de recherche que je lui avais proposé, savoir une théorie de structure des (Gr)-catégories, et la construction d\deleted{\unsure{b}}'une (Gr)-catégorie (ou plutôt catégorie de Picard) enveloppante pour une catégorie avec produit $\unsure{\otimes}$ associatif et commutatif. J'ai reçu récemment une copie du manuscript en forme (dont la confection s'est faite au prix de mille difficultés - il est douteux qu'il soit possible d'avoir d'autres copies dans un avenir proche, à moins d'en faire d'après cet exemplaire-là), et j'ai envoyé un rapport officiel au sujet de ce travail à M. Ta Quang B$\overset{\handwritten{\unsure{\tikz[baseline=0pt,x=.10em,y=.13em]{\draw[line width=.3pt] (0,0) .. controls (1.8,.2) and (2.0,1.6) .. (1.1,1.8);}|\,)}}}{\mathrm{u}}\overset{\handwritten{\unsure{\tikz[baseline=0pt,x=.14em,y=.11em]{\draw[line width=.3pt] (0,.85) .. controls (1.6,0) and (2.8,0) .. (3,.95) .. controls (3.05,1.0) and (3.1,1.1) .. (3.1,1.2);}|\,\sim}}}{\mathrm{u}}$, ministre de l'enseignement supérieur et technique de la RDV, rapport dont je vous envoie une copie ci-joint. Mme Sinh, avec mon accord, a demandé à présenter ce travail comme thèse de doctorat, et à faire sa soutenance en France. \deleted{\unsure{I|T}}A l'heure où j'écris, je n'ai pas connaissance d'un accord des autorités vietnamiennes pour l'un ni pour l'autre; il y a eu un certain nombre de difficultés soulevées, qu'il est intuile de détailler ici, et qui j'espère vont s'aplanir. J'ai appris par ailleurs que Mme Hoang \deleted{\unsure{H|V}}Xuan Sinh figurerait sur la liste officielle de la délégation vietnamienne au congrès international de mathématiciens à Vancouver du 20 août 1974. Si c'est bien ainsi, il serait très souhaitable qu'elle puisse profiter de cette occasion pour faire sa soutenance à Orsay au retour de Vancouver, fin Septembre ou début Octobre, et il me semble plausible que les autorités vietnamiennes y donneront leur accord. Aussi il serait assez urgent, vu les lenteurs administratives, de constituer un jury. Par la présente, je vous demande si vous seriez disposé en principe à faire partie (avec moi) de ce jury de thèse - ce qui suppose évidement que vous vous trouvez dans la région parisienne au moment prévu pour la soutenance. A celui qui m'en ferait la demande, je peux envoyer le texte de la thèse complète, ou un résumé assez détaillé rédigé par Mme Hoang Xuan Sinh , ou enfin un résumé moins détaillé de 13 pages manuscrites (rédigées en vue d'une publication aux CR, mais beaucoup trop long à cet effet). Ce serait gentil de me répondre rapidement, et en cas d'impossibilité, de voir autour de vous s'il y a des gens ``dans le bain'' d\deleted{e}\handwritten{u} \deleted{\unsure{ce}}\handwritten{genre} \handwritten{\inserted{de choses}} que fait Mme Sinh, et qui seraient intéressés à faire partie du jury.
% archive sheet number: 16
\scanpage{18}

Pour ce qui est des formalités administratives, c'est le frère de Mme Hoan Xuan \unsure{Han|Man}, \unsure{qui|qi} habite à Antony, qui s'en occupera pour elle. A toutes fin utile, je vous passe son adresse:

Hoang Xuan Sanh, 49 rue de Châtenay, Estérel, 92 Antony, Tél BER 68 79.\\

Dans l'atten\unsure{\uline{t}|t}e d'une réponse prochaine, bien cordialement\\
\hfill\handwritten{Schurik}\\

PS N'ayant pas l'adresse de Giraud, je demande à Verdier s'il peut bien lui transmettre la lettre et le rapport. Je\unsure{,}pense que celui-ci\unsure{\_}doit pouvoi\rlap{\raisebox{-1.1ex}[0pt][0pt]{\hspace{.2em}\unsure{(}}}r servir comme rapport de thèse aussi vis à vis de l'administration universitaire en France.
% archive sheet number: 17
